\subsection{Affine-proof interfaces, trace coordinates, and privacy limits}\label{sec:packing:affine}

\begin{definition}[Available affine constraints]\label{def:pack:affine-API}
An \emph{extension-field affine proof interface}
for $L$ allows a verifier to specify public
coefficients $(a_i)_i\in L^t$ and a public
right-hand side $b\in L$, and to require
committed hidden values $(x_i)_i$ to satisfy
\[
 \sum_i a_ix_i=b
\]
without exposing those values.
The algebraic interface does not itself
assert binding, knowledge soundness, or
zero knowledge for the implementation.
\end{definition}

\begin{lemma}[Incoming constraints are extension-field affine]\label{lem:pack:IN-affine}
For public $m,y',q_v,y_{\mathrm{pub}}$
the relations $\mathsf{CLAIM}$
and $\mathsf{IN}$ are extension-field
affine equations in hidden
$(p_v)_v,p_{\mathrm{in}}$.
\end{lemma}
\begin{proof}
The claim relation is
$p_{\mathrm{in}}=y'+y_{\mathrm{pub}}$
in characteristic two, so it is affine.
The incoming relation is
$\sum_vq_vp_v+p_{\mathrm{in}}
 =\sum_vq_vm_v+y'$,
with public right side and public
coefficients $q_v$. Both have exactly
the form in Definition~\ref{def:pack:affine-API}.
\end{proof}

\begin{lemma}[Outgoing constraints are extension-field affine]\label{lem:pack:OUT-affine}
For public $\theta,m$ and weights
$h_j$, the outgoing equation is
extension-field affine in hidden
$(\widehat w_j)_j$ and $c=(c_u)_u$.
\end{lemma}
\begin{proof}
Rewrite~\eqref{eq:pack:out-new} as
$\sum_ja_j(\theta)\widehat w_j+
  \sum_u\operatorname{eq}_\ell(u,\theta)c_u
  =B_m(\theta)$.
Every coefficient is publicly
computable from $\theta,h_j$,
and the public right side
$B_m(\theta)$ is computed by
transposing the public vector $m$.
No coordinate extraction of
a hidden field element occurs
in this equation, proving affinity.
\end{proof}

\begin{definition}[Affine correction coefficients]\label{def:pack:affine-H}
Use the trace-dual basis $\eta_u$ and
the public coefficients from Chapter~6:
\[
 H_{v,k}(\theta)
 =\sum_{u\in I}
    d_k\bigl(\beta_v\operatorname{eq}_\ell(u,\theta)\bigr)
    \eta_u.
\]
Define
\[
 V_k(p,c,\theta)
 =\eta_kD_c(\theta)+
       \sum_{v\in I}H_{v,k}(\theta)p_v.
\]
\end{definition}

\begin{lemma}[Every correction expression is extension-field affine]\label{lem:pack:V-affine}
For fixed public $\theta$,
the expression $V_k(p,c,\theta)$
is $L$-linear in the hidden values
$(p_v)_v$ and $(c_u)_u$.
\end{lemma}
\begin{proof}
Substitute
$D_c(\theta)=\sum_u\lambda_u(\theta)c_u$.
Then
\[
 V_k=\sum_u\eta_k\lambda_u(\theta)c_u
       +\sum_vH_{v,k}(\theta)p_v.
\]
All coefficients multiplying hidden
variables depend only on public
parameters and $\theta$.
There is no constant term, so the
expression is in fact homogeneous
$L$-linear.
\end{proof}

\begin{theorem}[Trace-coordinate interface for candidate correction]\label{thm:pack:trace-correct}
For every $p,c,\theta$,
\[
 \boxed{
 D_c(\theta)=\Delta_p(\theta)
 \quad\Longleftrightarrow\quad
 \Tr(V_k(p,c,\theta))=0\quad\forall k\in I.}
\]
\end{theorem}
\begin{proof}
By Theorem~\ref{thm:binary:dual-basis},
$d_k(x)=\Tr(\eta_kx)$.
Lemma~\ref{lem:binary:affine-trace}
shows
$\Tr(\sum_vH_{v,k}p_v)
 =d_k(\Delta_p)$.
Therefore
$\Tr(V_k)=d_k(D_c+\Delta_p)$,
which is Lemma~\ref{lem:binary:V-trace}.
Every coordinate vanishes if and only
if the field element vanishes,
by the basis expansion.
In characteristic two
$D_c+\Delta_p=0$ is equivalent to
$D_c=\Delta_p$.
\end{proof}

\begin{lemma}[A revealed transposed pad reveals the original pad]\label{lem:pack:reveal-pad}
If the entire field vector $c^p=\mathcal T(p)$
is revealed, then the pad vector $p$
is uniquely recovered by
\[
 \boxed{p=\mathcal T(c^p).}
\]
\end{lemma}
\begin{proof}
Apply $\mathcal T$ to the identity
$c^p=\mathcal T(p)$, then use
Theorem~\ref{thm:pack:transpose-involutive}.
\end{proof}

\begin{theorem}[Revealing the correction destroys masking]\label{thm:pack:reveal-leak}
Suppose both $m=s^*+p$ and
$c^p=\mathcal T(p)$ are sent in the clear.
Then every unmasked partial evaluation
$s_v^*$ is deterministically recoverable:
\[
 \boxed{s^*=m+\mathcal T(c^p).}
\]
\end{theorem}
\begin{proof}
Lemma~\ref{lem:pack:reveal-pad}
recovers $p=\mathcal T(c^p)$.
As $m=s^*+p$, adding $p$ to
both sides yields
$s^*=m+p=m+\mathcal T(c^p)$.
Thus the unmasked vector can be
computed from the two exposed vectors.
\end{proof}

\begin{corollary}[Masked-value uniformity alone does not imply privacy]\label{cor:pack:uniform-not-ZK}
The marginal uniformity result
for $m$ in Lemma~\ref{lem:pack:masked-uniform}
does not imply witness privacy for the
joint message pair $(m,c^p)$.
\end{corollary}
\begin{proof}
For any fixed $s^*$, the joint pair
$(m,c^p)$ deterministically reveals
$s^*$ by Theorem~\ref{thm:pack:reveal-leak}.
Choose two distinct possible
partial-evaluation vectors $s_0^*,s_1^*$.
The events that the reconstruction
equals $s_0^*$ distinguish their joint
message distributions with probabilities
$1$ and $0$, respectively.
Thus the pair need not be
witness independent, despite each
$m$ being uniform.
\end{proof}

\begin{remark}[Public parameters and hidden values]
All $H_{v,k}(\theta)$ and $a_j(\theta)$
may be computed by the verifier from
public inputs; the vectors $p,c$
must remain hidden behind a commitment
or an appropriate zero-knowledge
affine proof. Whether such a proof can
simulate its commitments, openings and
messages is a separate cryptographic
statement. The current chapter proves
only the algebraic affine interface.
\end{remark}

\begin{remark}[Independent paper contributions]
The matrix transposition and packed
bit-linear reduction originate in the
Ring Switching line of work
\cite{DiamondPosen2024}.
The trace/Frobenius approach to hiding
has prior treatments
\cite{Kochergin2026,FrigoShelat2026}.
Our kernel-coordinate identities use
the basis from KBFold
\cite{MkhidaKBFold2026}.
No novelty claim about the
underlying trace-zero principle is
made here.
\end{remark}
